\begin{sloppypar}

\subsection{Fenster}

Die Klasse \texttt{czg.MainWindow} stellt das Spielfenster
dar und beinhaltet die \texttt{main}-Methode sowie eine
Reihe von statischen Konstanten. Letztere werden in der
Dokumentation mit dem gleichen Namen wie im Code bezeichnet.

\begin{itemize}
	\item \texttt{int PIXEL\_SCALE}: Da das Spiel im Pixel-Art-Stil
	      umgesetzt wurde, sind die im Spiel enthaltenen Bilddateien sehr
	      klein. Deshalb werden sie um diesen konstanten Faktor skaliert.
	\item \texttt{int WIDTH, HEIGHT}: Die Höhe und Breite des Spielfensters
	      in Pixeln.
	\item \texttt{int FPS}: Abkürzung für \textit{Frames per Second}.
	      Wie oft Spiellogik und Rendern in einer Sekunde ausgeführt
	      werden sollen.
	\item \texttt{MainWindow INSTANCE = new MainWindow()}: Die \textit{Singleton}-
	      Instanz des Fensters.
	\item \texttt{long TIME\_AT\_UPDATE\_START}: Der Wert von \texttt{System.nanoTime()}
	      zu beginn eines Logik-Durchlaufs. Siehe Abschnitt \ref{programm:engine:eingabe}.
\end{itemize}


\subsubsection{\texttt{main}-Methode}

\texttt{czg.MainWindow.main} erledigt folgende Initialisierungsaufgaben:

\begin{itemize}
	\item Platzieren das Fenster in der Mitte des Bildschirms und
	      Anzeigen desselben
	\item Starten der Hintergrundmusik, indem eine neue Instanz von
	      \texttt{czg.sound.StreamSound} zur globalen Sound-Gruppe
	      \texttt{czg.sound.SoundGroup.GLOBAL\_SOUNDS} hinzugefügt wird.
	      Siehe Abschnitt \ref{programm:engine:sound:gruppe}.
	\item Korrigieren der Fenstergröße. Dies ist notwendig, da beim
	      Aufruf von \texttt{javax.swing.JFrame.setSize} die vom
	      Betriebssystem angezeigte Fensterleiste mit einberechnet
	      ist. Entsprechend muss die Höhe des Fensters noch einmal
	      um die Höhe der Fensterleiste erhöht werden, um keinen Teil
	      der Grafik abzuschneiden.
	\item Anzeigen der ersten Szene, indem eine neue Instanz der entsprechenden
	      Klasse zum Szenenstapel hinzugefügt wird. Siehe Abschnitt \ref{programm:engine:szene:stapel}.
	\item Starten der Hauptschleife. Siehe Abschnitt \ref{programm:engine:hauptschleife}.
\end{itemize}


\subsubsection{Hauptschleife} \label{programm:engine:hauptschleife}

Wie in Abschnitt \ref{idee:engine:hauptschleife} erwähnt läuft die
Hauptschleife mit einer Frequenz von 60Hz, wie in \texttt{czg.MainWindow.FPS}
definiert. Die Grundidee des Algorithmus findet sich ebenfalls dort.
Im folgenden ist die letztendliche Implementierung in \texttt{MainWindow.java}
gelistet:

\begin{lstlisting}[caption=Hauptschleife, label=programm:engine:hauptschleife:code, captionpos=t, language=Java]
// In welchem Zeitintervall (in Nanosekunden) die Spiellogik
// ausgeführt werden soll
final double interval = 1e9 / FPS;
// Zählt, wie oft die Spiellogik durchlaufen werden sollte.
double delta = 0;

// Zeit seit dem letzten Durchlauf der while(true)-Schleife
long lastTime = System.nanoTime();
// Zweite Zeit-Variable. Wird zur neuen lastTime.
long currentTime;

//noinspection InfiniteLoopStatement
while(true) {
	// Aktuelle Zeit messen
	currentTime = System.nanoTime();
	
	// Die Zeit berechnen, die seit dem letzten Durchlauf
	// dieser Schleife vergangen ist.
	// Damit berechnen, wie oft die Spiellogik in dieser
	// Zeitspanne hätte durchlaufen
	// werden sollen (normalerweise eine Anzahl <1).
	delta += (currentTime - lastTime) / interval;
	
	// Die currentTime wird wieder zur lastTime
	lastTime = currentTime;
	
	// Alle nötigen Durchläufe abarbeiten
	while(delta >= 1) {
		TIME_AT_UPDATE_START = System.nanoTime();
		
		// Code für Szenen und Objekte ausführen
		SceneStack.INSTANCE.update();
		
		[ausgelassen]
		
		// Grafik
		SceneStack.INSTANCE.repaint();
		
		// Zusätzliche Aktualisierung
		// für die Eingabeverarbeitung
		Input.INSTANCE.update();
		
		// Durchlauf abgeschlossen, Zähler kann um
		// 1 verringert werden
		delta--;
	}
}
\end{lstlisting}

Der hier ausgelassene Teil des Codes beinhaltet lediglich
die Logik für einige Tastenkürzel. Diese werden in
späteren Abschnitten an den entsprechenden Stellen
erwähnt und sind im Programmablaufprotokoll noch einmal
gelistet.


\subsection{Eingabe} \label{programm:engine:eingabe}

Die Klasse \texttt{czg.util.Input} implementiert die in
\ref{idee:engine:grafik-eingabe} erwähnten Interfaces
\texttt{java.awt.event.KeyListener} und \texttt{java.awt.event.MouseListener},
sowie zusätzlich noch \texttt{java.awt.event.FocusListener}.
Zusätzlich definiert sie die Konstante \texttt{HELD\_THRESHOLD}.
Diese sagt aus, wie lange eine Taste gedrückt sein darf und
dabei noch als gedrückt und nicht als gehalten gilt.
Die Verwendung dieser Konstante sowie die verschiedenen
Zustände von Tasten werden im Verlauf dieses Abschnitts
näher erklärt.

Der Konstruktor der Klasse ist privat. Es kann jedoch
über die statische Konstante \texttt{Input.INSTANCE}
auf eine Singleton-Instanz der Klasse zugegriffen werden.
Diese wird u.A. von \texttt{MainWindow} an deren
\texttt{addKeyListener}- und \texttt{addMouseListener}-Methoden
übergeben, um das Eingabeverarbeitungssystem funktionstüchtig
zu machen.


\subsubsection{Speichern von gedrückten Tasten}

Den Kern der \texttt{Input}-Klasse bilden zwei Instanzen von
\texttt{java.util.concurrent.ConcurrentHashMap}:
\texttt{keyStates} und \texttt{mouseStates}.
Wird eine Taste gedrückt, werden die entsprechenden
Methoden \texttt{keyPressed} und \texttt{mousePressed} von
\texttt{Swing} aufgerufen. Die Implementierungen dieser
Methoden in \texttt{Input} ordnen dann dem Tasten-Code der
gedrückten Taste, welcher über die \texttt{getKeyCode}-Methode
des \texttt{KeyEvent}- bzw. \texttt{MouseEvent}-Objekts,
welches an die Methode übergeben wird, erhalten werden kann,
den aktuellen Wert von \texttt{System.nanoTime()} zu, speichern
also den Zeitpunkt, an welchem die Taste gedrückt wurde.

Analog wird ein Eintrag in der entsprechenden \texttt{ConcurrentHashMap}
wieder entfernt, wenn die zugehörige Taste losgelassen und
somit \texttt{Input.keyReleased} bzw. \texttt{Input.mouseReleased}
aufgerufen wird.

Befindet sich das Spielfenster nicht mehr im Vordergrund bzw. im Fokus,
so ruft \texttt{Swing} die Methode \texttt{focusLost} aus dem \texttt{FocusListener}-Interface
auf. Die Implementierung dieser Methode in \texttt{Input} besteht
daraus, alle Einträge aus den beiden \texttt{ConcurrentHashMap}s zu
entfernen.


\subsubsection{Tastenzustände und das Abfragen dieser}

Nun kann der Zustand einer Taste mittels den Methoden
\texttt{Input.getKeyState(int)} und \texttt{Input.getMouseState(int)}
abgefragt werden. Die zu übergebenden Parameter sind die
Tasten-Codes der abzufragenden Maus- oder Tastaturtaste,
welche als statische Konstanten in den bereits erwähnten
Klassen \texttt{KeyEvent} und \texttt{MouseEvent} definiert sind.

Die beiden Methoden geben eine Konstante aus dem Enum
\texttt{Input.KeyState} zurück. Dieses definiert genau
drei Zustände für Tasten:

\begin{itemize}
	\item \texttt{NOT\_PRESSED}: Die Taste ist nicht gedrückt
	\item \texttt{PRESSED}: Die Taste wurde, von \texttt{MainWindow.TIME\_AT\_UPDATE\_START}
	      aus gemessen, höchstens für die Länge des Zeitintervalls
	      \texttt{Input.HELD\_THRESHOLD} gedrückt.
	\item \texttt{HELD}: Ist eine Taste schon länger gedrückt
	      als \texttt{Input.HELD\_THRESHOLD}, gilt sie als gehalten
	      und nicht mehr als gedrückt.
\end{itemize}

Zusätzlich bietet \texttt{KeyState} die Methode \texttt{isDown},
welche \texttt{true}, wenn es sich bei der Enum-Konstante,
für welche die Funktion aufgerufen wurde, um \texttt{PRESSED}
oder \texttt{HELD} handelt, und sonst \texttt{false} zurückgibt.

Die von \texttt{Input.getKeyState(int)} verwendete
statische Methode \texttt{KeyState.fromTimePressed(long)} verlangt
als Parameter den Zeitpunkt, an dem die Taste gedrückt wurde, also
\texttt{keyPressed} bzw. \texttt{mousePressed} aufgerufen wurde,
leitet daraus die passende Enum-Konstante ab und gibt sie zurück.

Wie im Programmcodelisting \ref{programm:engine:hauptschleife:code} der Hauptschleife
zu sehen ist, wird als letzter Teil jedes Frames die Methode
\textit{Input.update} aufgerufen. Dies hängt zusammen mit den
internen Listen \texttt{Input.keyButtonsToUpdateToHeld} und
\texttt{Input.mouseButtonsToUpdateToHeld}: Jedes mal, wenn
\texttt{Input.getKeyState} bzw. \texttt{Input.getMouseState}
feststellt, dass der Zustand einer Taste \texttt{PRESSED} ist,
wird deren Tasten-Code zu der entsprechenden Liste hinzugefügt.
Beim Aufruf von \texttt{Input.update} werden dann alle Einträge
für die Tasten-Codes in den Listen in der entsprechenden Map
zu \texttt{HELD} geändert - wurde eine Taste also in einem Frame
abgefragt und galt als \texttt{PRESSED}, so gilt sie ab dem nächsten
Frame immer als \texttt{HELD}, egal wie lange sie eigentlich gehalten
wurde. So werden z.B. Mausklicks nicht über mehrere Frames hinweg
mehrfach registriert, denn die entsprechende Maustaste kann nur
für höchstens einen Frame als \texttt{PRESSED} gelten.


\subsubsection{Mausposition}

Außerdem speichert \texttt{Input} die Position des Mauszeigers
sowie diese im vorangegangenen Frame. Dies ist gelöst mithilfe
zweier privater Variablen \texttt{mousePosition} und
\texttt{lastMousePosition}, deren Werte über entsprechende
Getter-Methoden abgefragt werden können. Beim Aufruf von
\texttt{Input.update} erhält \texttt{lastMousePosition} den
Wert von \texttt{mousePosition}. Letztere bekommt dann einen
neuen Wert, welcher mit der \texttt{getMousePosition}-Methode
des Szenenstapels erhalten wird, da dieser \texttt{javax.swing.JPanel}
erweitert. Siehe \ref{programm:engine:szene:stapel}. Würde
stattdessen die Methode mit dem gleichen Namen von \texttt{MainWindow}
aufgerufen werden, so wäre die Mausposition versetzt, da der
Ursprung der von dieser Methode zurückgegebenen Koordinaten
bei der oberen linken Ecke des Fensters \textbf{inklusive Titelleiste}
liegt.

\subsection{Spiel-Objekte}

Spiel-Objekte werden mithilfe der Klasse \texttt{czg.objects.BaseObject}
umgesetzt. Diese speichert folgende Eigenschaften:

\begin{itemize}
	\item X- und Y-Koordinaten, in Pixeln. Dabei ist $(0|0)$ die
		  obere linke Ecke des Spielfensters, und die Koordinaten
		  steigen nach unten bzw. nach rechts in positive Richtung an.
	\item Höhe sowie Breite des Objekts, in Pixeln
	\item Ein Bild, welches an den Koordinaten und mit der Größe
	      des Objekts gezeichnet werden soll. Dieses kann auch
	      \texttt{null} sein, in diesem Fall ist das Objekt
	      unsichtbar.
\end{itemize}

Diese sind alle als \texttt{public} deklariert, können also
von jeder Stelle im Code aus gelesen und geändert werden.

Die Klasse bietet drei Konstruktoren, welche verschiedene
Teilmengen der oben gelisteten Eigenschaften als Parameter
annehmen: Entweder nur ein Bild, ein Bild und eine Position,
oder alle drei Eigenschaften. Werden Höhe und Breite nicht
explizit angegeben, werden sie auf die Höhe und Breite des
Bildes, jeweils mit \texttt{PIXEL\_SCALE} multipliziert,
gesetzt. Wird keine Position angegeben, wird das Objekt
in die Mitte des Spielfensters platziert.

Weiterhin definiert \texttt{BaseObject} eine Reihe an
Methoden:

\begin{itemize}
	\item \texttt{getHitbox}: Verpackt Position und Größe in ein
	      \texttt{java.awt.geom.Rectangle2D}-Objekt, welches eine
	      Reihe nützlicher Methoden bietet.
	\item \texttt{isClicked(boolean includeTransparency)}: Prüft,
	      ob das Objekt angeklickt wurde, also ob sich der Mauszeiger
	      über dem Objekt befindet und die linke Maustaste gedrückt
	      ist. Der optionale Parameter \texttt{includeTransparency}
	      gibt an, ob auch ein Klick auf transparente Bereiche
	      des Bildes des Objektes berücksichtigt werden soll.
	\item \texttt{draw(Graphics2D)}: Die vorgegebene Implementierung
	      dieser Methode zeichnet das Bild des Objektes mit der
	      gespeicherten Breite und Höhe an der gespeicherten Position
	      mithilfe der \texttt{drawImage}-Methode des an die Methode
	      übergebenen \texttt{Graphics2D}-Objektes. Letzteres stammt
	      von der \texttt{paintComponent}-Methodes des Szenenstapels.
	      Siehe \ref{programm:engine:szene:stapel}.
	\item \texttt{update(BaseScene)}: Wird vom Szenen-Stapel bzw.
	      der Szene, in welcher sich das Objekt befindet, aufgerufen.
	      Letztere ist dieselbe, welche an die Methode als Parameter
	      übergeben wird. Die Implementierung dieser Methode in
	      \texttt{BaseObject} ist leer und kann von Unterklassen
	      überschrieben werden, die eigene Logik implementieren wollen.
	\item \texttt{unload(BaseScene)}: Wird aufgerufen, wenn die Szene,
	      in welcher sich das Objekt befindet, entfernt wird. Letztere
	      ist, analog zu \texttt{update}, ebenfalls dieselbe, in der
	      sich das Objekt befindet. Auch hier ist die Implementierung in
	      \texttt{BaseObject} leer und kann bei bedarf von Unterklassen
	      überschrieben werden.
\end{itemize}

Sowohl \texttt{update} als auch \texttt{unload} waren ursprünglich als
\texttt{abstract} markiert. Jedoch brachte dies den Nachteil, dass
\texttt{BaseObject} somit ebenfalls abstrakt wurde und somit nicht
mehr direkt instanziiert werden konnte. Jedoch wurden an einigen
Stellen im Code Objekte benötigt, welche keine weitere Funktion als
das darstellen eines Bildes haben sollten, weshalb es also wieder
ermöglicht werden musste, \glqq{}leere\grqq{} \texttt{BaseObject}s
ohne weiteren Code zu erstellen.


\subsection{Szenen}

Eine Szene ist eine Gruppierung von Objekten. Die Umsetzung erfolgt
über die abstrakte Klasse \texttt{czg.scenes.BaseScene}. Diese speichert
folgende Daten:

\begin{itemize}
	\item Eine Liste von \texttt{BaseObject}s. Diese ist zwar \texttt{final},
	      aber auch \texttt{public}. Es können also von jeder Stelle beliebig
	      Objekte hinzugefügt und entfernt werden.
	\item Eine eigene Sound-Gruppe. Siehe \ref{programm:engine:sound:gruppe}.
	\item Eine Liste von \textit{Tags}. Siehe \ref{programm:engine:szene:verdecken}.
\end{itemize}

Daneben besitzt \texttt{BaseScene} folgende wesentliche Methoden:

\begin{itemize}
	\item \texttt{draw(Graphics2D)}: Das \texttt{Graphics2D}-Objekt
	      stammt vom Szenenstapel. Siehe \ref{programm:engine:szene:stapel}.
	      Die vorgegebene Implementierung dieser Methode iteriert
	      einfach über die Liste der Objekte in der Szene und ruft deren
	      \texttt{draw(Graphics2D)}-Methode auf, wobei das an die
	      \texttt{draw}-Methode übergebene \texttt{Graphics2D}-Objekt
	      direkt an die gleichnamige Methode der Objekte übergeben wird.
	\item \texttt{update(BaseScene)}: Wird vom Szenen-Stapel aufgerufen.
		  Es wird eine Kopie der Objekt-Liste angelegt, über welche dann
		  iteriert wird. Dabei wird die \texttt{update(BaseScene)}-Methode
		  jedes Objekts aufgerufen und die Szene selbst als Parameter
		  übergeben. Der Grund dafür, dass über eine Kopie der Liste
		  iteriert wird, ist, dass die Logik der Objekte das ändern
		  der Objekteliste ihrer Szene beinhalten könnte. Das Ändern
		  einer Liste während über diese iteriert wird resultiert
		  jedoch in einer \texttt{ConcurrentModificationException}.
		  Diese wird so umgangen.
	\item \texttt{unload(BaseScene)}: Wird aufgerufen, wenn diese
	      Szene vom Szenenstapel entfernt wird. Die sich in der Sound-Gruppe
	      der Szene befindenden Sounds werden gestoppt und die
	      \texttt{unload}-Methode aller Objekte in der Szene
	      wird aufgerufen.
	      Siehe auch \ref{programm:engine:szene:stapel}.
\end{itemize}

Diese drei Methoden sollten von Unterklassen bei bedarf überschrieben
werden, wobei aber dennoch ein Aufruf der entsprechenden Methode,
wie sie in \texttt{BaseScene} implementiert ist, stattfinden sollte.

Die Konstruktoren sowie einige weitere Variablen und Methoden
der Klasse sind in den kommenden zwei Abschnitten erklärt.


\subsection{Der Szenenstapel} \label{programm:engine:szene:stapel}

Um mehrere Szenen übereinander anzeigen und deren Logik gemeinsam
ausführen zu können, steht der \textit{Szenenstapel} in Form der
Klasse \texttt{czg.scenes.SceneStack} bzw. deren Singleton-Instanz
\texttt{SceneStack.INSTANCE} zur Verfügung. Diese erweitert
\texttt{javax.swing.JPanel} und ist die einzige Komponente, die sich
im Spielfenster befindet.

Die Klasse speichert die folgenden Daten, welche beide als \texttt{private}
deklariert sind:

\begin{itemize}
	\item Die Liste der Szenen auf dem Stapel
	\item Eine weitere Liste, welche für jede Szene im Stapel
	      die Vereinigungsmenge der Tags der darüber liegenden
	      Szenen speichert und wie oft jeder Tag jeweils darüber
	      liegt. Siehe \ref{programm:engine:szene:verdecken}.
\end{itemize}

Jede \texttt{BaseScene} speichert außerdem ihren Index in der
Liste des Szenenstapels. Ist die Szene nicht im Stapel, wird
dieser auf $-1$ gesetzt.


\subsubsection{Operationen}

Um den Szenenstapel zu ändern, werden folgende als
\texttt{public} deklarierte Methoden bereitgestellt:

\begin{itemize}
	\item \texttt{push(BaseScene)}: Fügt die gegebene Szene ganz oben auf
	      dem Stapel, also am Ende der Liste ein.
	\item \texttt{pop}: Entfernt die Szene, die ganz oben auf dem Stapel
	      liegt, also jene, die zuletzt hinzugefügt wurde.
	\item \texttt{insert(BaseScene, int)}: Fügt die gegebene Szene an dem
	      gegebenen Index im Stapel ein.
	\item \texttt{replace(BaseScene, BaseScene)}: Sucht die erste gegebene
	      Szene im Stapel und ersetzt sie durch die zweite gegebene Szene.
	\item \texttt{remove(BaseScene)}: Entfernt die gegebene Szene.
	\item \texttt{clear}: Entfernt alle Szenen aus dem Stapel.
\end{itemize}

Beim Entfernen einer Szene wird deren \texttt{unload}-Methode aufgerufen.


\subsubsection{Grafik}

Um die Szenen im Stapel und deren Objekte zu zeichnen, überschreibt
\texttt{SceneStack} zunächst die \texttt{paintComponent(Graphics)}-Methode
von \texttt{JPanel}. Von dort wird dann die eigentliche Methode aufgerufen,
die für die Grafik verantwortlich ist: \texttt{draw(Graphics2D)}. Diese
übernimmt das nun zu \texttt{Graphics2D} umgewandelte \texttt{Graphics}-Objekt,
welches von \texttt{Swing} für \texttt{paintComponent} bereitgestellt wird.
\texttt{SceneStack.draw} aktiviert dann über \texttt{Graphics2D.setRenderingHint}
die \textit{Anti-Aliasing}-Funktion. Diese sorgt dafür, dass vor allem Text
nicht verpixelt und scharfkantig, sondern ansehnlich und mit weichen
Konturen dargestellt wird. Dann delegiert \texttt{draw} das eigentliche
Rendern an die \texttt{SceneStack.processScenes}-Methode. Diese wird
in kürze erklärt.

Das Auslagern des Grafikcodes in eine separate Methode ist nötig, damit
diese auch ohne eine grafische Umgebung aufgerufen werden kann, also von
Tests aus, die z.B. automatisch von GitHub ausgeführt werden, wenn neue
Commits hochgeladen wurden. In diesem Fall übergibt der Test-Code \texttt{null}
an \texttt{SceneStack.draw}. Dann wird ein neuer Grafikkontext erstellt,
welcher nicht an ein Fenster gebunden ist, sondern nur im Speicher existiert,
aber dennoch von Szenen und Objekten verwendet werden kann.


\subsubsection{Logik}

Das Ausführen der Logik von Szenen und Objekten wird
über die Methode \texttt{SceneStack.update} erreicht.
Auch diese delegiert an die bereits erwähnte
 \texttt{processScenes}-Methode.


\subsubsection{Iteration und Verarbeitung der Szenen}

Die Methode \texttt{SceneStack.processScenes} erhält drei Parameter:

\begin{enumerate}
	\item Eine \texttt{java.util.function.Function}, die ein Verdeckungs-Einstellungen-Objekt
	      annimmt und die spezifische Einstellung, welche bestimmt, ob die Szene bei
	      der Iteration über die Szenenliste einbezogen werden soll, extrahiert.
	      \begin{itemize}
	          \item \texttt{update}: Referenz zur Getter-Methode für \texttt{coverPausesLogic}
	          \item \texttt{draw}: Referenz zur Getter-Methode für \texttt{coverDisablesDrawing}
	      \end{itemize}
	      Siehe \ref{programm:engine:szene:verdecken}.
	\item Einen \texttt{java.util.Consumer<BaseScene>}, welcher bestimmt, wie mit
	      den einzelnen Szenen bei der Iteration zu verfahren ist.
 	      \begin{itemize}
	      	\item \texttt{update}: Referenz zu \texttt{BaseObject.update}
	      	\item \texttt{draw}: Referenz zu \texttt{BaseObject.draw}
	      \end{itemize}
	\item Einen \texttt{boolean}, der angibt, ob eine Kopie der Szenenliste
	      für die Iteration verwendet werden sollte oder nicht.
\end{enumerate}

Zunächst werden nun, wenn angegeben, jeweils eine Kopie der Szenenliste und der
Liste der über jeder Szene liegenden Tags erstellt. Dann wird über die Szenen
iteriert:

\begin{enumerate}
	\item Anhand der Tags der über der Szene liegenden Szenen wird
	      die effektiven Verdeckungs-Einstellungen bestimmt. Dann wird
	      mit der als ersten Parameter übergebenen Funktion die entsprechende
	      Einstellung daraus extrahiert.
	\item \begin{enumerate}
		\item Ist Szene Verdeckt und die Einstellung auf \textit{aktiv} gestellt,
		      wird die Szene übersprungen und es wird mit der nächsten fortgefahren.
		\item Andernfalls wird der als zweiter Parameter angegebene \texttt{Consumer}
		      auf die Szene angewendet.
	\end{enumerate}
\end{enumerate}


\subsection{Verdecken von Szenen} \label{programm:engine:szene:verdecken}

\texttt{BaseScene} bietet zwei Konstruktoren, einen ohne Parameter und einen
mit den Verdeckungs-Einstellungen als einzigen Parameter. Ersterer ruft lediglich
letzteren auf und übergibt ein \texttt{CoverSettings}-Objekt mit Standardeinstellungen.
Siehe \ref{programm:engine:szene:verdecken}.

\begin{itemize}
	\item Eine \textit{boolean}-Variable, die speichert ob die Szene von einer
	anderen verdeckt wird, also ob im Szenen-Stapel eine andere über ihr liegt.
	Siehe \ref{programm:engine:szene:stapel} und \ref{programm:engine:szene:verdecken}.
	\item Die \textit{Verdeckungs-Einstellungen}, die angeben, was passieren soll,
	wenn die Szene von anderen Szenen, ggf. solche mit bestimmten Tags, verdeckt ist.
	Siehe \ref{programm:engine:szene:stapel} und \ref{programm:engine:szene:verdecken}.
\end{itemize}

\subsection{Sound}

Die \texttt{Java Sound API} untersetützt standardmäßig nur das unkomprimierte \texttt{WAV}-Format. Die Audiodateien liegen allerdings 
aus Gründen der effizienteren Speicherplatznutzung im komprimierten \texttt{Ogg/Vorbis}-Format vor. Um diese laden zu können

\subsection{Sound-Gruppe} \label{programm:engine:sound:gruppe}

\end{sloppypar}